\section*{Chapitre \ref{ch:b1:elementary-steps} --- Mécanismes réactionnels~: actes élémentaires et approximations}

\begin{solution}{exo:b1:elementary-steps:1}
(a) Monomoléculaire~: peut être élémentaire. (b) Bimoléculaire~: peut
être élémentaire. (c) Trois molécules et quatre liaisons reconstruites à
la fois~: pas élémentaire. (d) Trimoléculaire, \ce{N2} emportant
l'énergie libérée~: un \omterm{def:b1:elementary-steps:elementary-step}{acte élémentaire} rare mais réel.
\end{solution}

\begin{solution}{exo:b1:elementary-steps:2}
$v = k[\mathrm{A}][\mathrm{B}]$~; $v = k[\mathrm{A}]^2$~; $v = k[\mathrm{A}]$.
Une équation globale est un bilan, non une description de la façon dont
les molécules se rencontrent~; sa \omterm{def:b1:rate-laws:order}{loi de vitesse} doit être mesurée, ou
déduite d'un mécanisme.
\end{solution}

\begin{solution}{exo:b1:elementary-steps:3}
L'intermédiaire est le creux à 30~; les \omterm{def:b1:elementary-steps:energy-profile}{états de transition} sont les
maxima à 70 et 55. L'\omterm{def:b1:elementary-steps:energy-profile}{état de transition} le plus haut est le premier~: le
premier acte est cinétiquement déterminant. Il est endothermique (de 0
à 30).
\end{solution}

\begin{solution}{exo:b1:elementary-steps:4}
Par le même facteur $10^4$ (\cref{prop:b1:elementary-steps:catalysis})~;
la constante d'équilibre est inchangée~; l'équilibre est atteint environ
$10^4$ fois plus tôt.
\end{solution}

\begin{solution}{exo:b1:elementary-steps:5}
$t_{\max} = \ln(0.30/0.10)/(0.30 - 0.10) = \qty{5.49}{s}$. Alors
$[\mathrm{B}]_{\max} = a_0\frac{0.10}{0.20}(\eu^{-0.549} - \eu^{-1.648}) =
0.5\,a_0(0.577 - 0.192) = 0.192\,a_0$.
\end{solution}

\begin{solution}{exo:b1:elementary-steps:6}
$\dd[\mathrm{I}]/\dd t = k_1[\mathrm{A}] - (k_{-1} + k_2)[\mathrm{I}] = 0$
donne $[\mathrm{I}] = k_1[\mathrm{A}]/(k_{-1} + k_2)$ et $v = k_2[\mathrm{I}]
= \frac{k_1k_2}{k_{-1}+k_2}[\mathrm{A}]$. Si $k_2 \gg k_{-1}$, $v =
k_1[\mathrm{A}]$~: le premier acte est cinétiquement déterminant. Si
$k_2 \ll k_{-1}$,
$v = (k_1/k_{-1})k_2[\mathrm{A}]$~: un pré-équilibre suivi d'un acte
lent.
\end{solution}

\begin{solution}{exo:b1:elementary-steps:7}
Somme~: \ce{2H2O2 -> 2H2O + O2}. \omterm{def:b1:elementary-steps:catalyst}{Catalyseur}~: \ce{I-} (consommé, puis
régénéré)~; intermédiaire~: \ce{IO-}. Le premier acte, lent, fixe la
vitesse~: $v = k_1[\ce{H2O2}][\ce{I-}]$.
\end{solution}

\begin{solution}{exo:b1:elementary-steps:8}
L'acte est fortement endothermique~: d'après le postulat de Hammond, son
\omterm{def:b1:elementary-steps:energy-profile}{état de transition} ressemble au carbocation. Un substituant qui abaisse
l'énergie du carbocation abaisse presque autant celle de l'\omterm{def:b1:elementary-steps:energy-profile}{état de
transition}~: la barrière diminue et l'acte est plus rapide.
\end{solution}

\begin{solution}{exo:b1:elementary-steps:9}
$k_{\mathrm B}/k_{\mathrm C} = \exp(20000/RT)$~: \num{3.0e3} à
\qty{300}{K}, $\exp(4.01) = 55$ à \qty{600}{K}. Le chauffage réduit la
préférence cinétique pour B et, en rendant renversable la formation de
B, permet au système d'atteindre C, plus stable.
\end{solution}

\begin{solution}{exo:b1:elementary-steps:10}
$k_1[\mathrm{A}][\mathrm{M}] = (k_{-1}[\mathrm{M}] + k_2)[\mathrm{A}^*]$,
donc $v = k_2[\mathrm{A}^*] = \dfrac{k_1k_2[\mathrm{A}][\mathrm{M}]}{k_{-1}[\mathrm{M}]
+ k_2}$. À haute pression ($k_{-1}[\mathrm{M}] \gg k_2$), $v =
(k_1k_2/k_{-1})[\mathrm{A}]$~: ordre 1. À basse pression, $v =
k_1[\mathrm{A}][\mathrm{M}]$~: ordre 2, l'activation par choc devient
l'acte lent.
\end{solution}

\begin{solution}{exo:b1:elementary-steps:11}
La démonstration est celle du théorème. Pour $k_1 = k_2 = k$~:
$\frac{\dd}{\dd t}([\mathrm{B}]\eu^{kt}) = ka_0$, donc $[\mathrm{B}]\eu^{kt} =
ka_0t$ et $[\mathrm{B}] = a_0kt\,\eu^{-kt}$, maximale en $t = 1/k$.
\end{solution}

\begin{solution}{exo:b1:elementary-steps:12}
$\mathrm{A} + \mathrm{B} \rightleftharpoons \mathrm{I} + \mathrm{C}$
(rapide, de constante $K_1$), puis $\mathrm{I} + \mathrm{B} \to \mathrm{P}$
(lent,
$k_2$). La somme est la réaction globale~; $[\mathrm{I}] =
K_1[\mathrm{A}][\mathrm{B}]/[\mathrm{C}]$ et $v = k_2[\mathrm{I}][\mathrm{B}]
= k_2K_1[\mathrm{A}][\mathrm{B}]^2/[\mathrm{C}]$~: le sous-produit, libéré
dans le pré-équilibre, ralentit la réaction.
\end{solution}

\begin{solution}{pb:b1:elementary-steps:1}
\textbf{1.} Doubler $[\ce{NO}]_0$ multiplie $v_0$ par 4~: ordre 2~;
doubler $[\ce{O2}]_0$ la multiplie par 2~: ordre 1.
\textbf{2.} $v = k[\ce{NO}]^2[\ce{O2}]$~; $k = 1.0 \times 10^{-5}/((1.0
\times 10^{-3})^2 \times 1.0 \times 10^{-3}) =
\qty{1.0e4}{L^2.mol^{-2}.s^{-1}}$.
\textbf{3.} Un choc simultané de trois molécules est rare, et un acte
unique ne pourrait pas devenir plus lent quand on chauffe (question 5).
\textbf{4.} $E_a = R\ln(k_{400}/k_{300})/(1/300 - 1/400) = 8.314 \ln(0.1)
/(8.33 \times 10^{-4}) = \qty{-23}{kJ/mol}$.
\textbf{5.} Un \omterm{def:b1:elementary-steps:elementary-step}{acte élémentaire} procède par des chocs qui franchissent
une barrière~; la fraction de ceux qui y parviennent, $\eu^{-E_a/RT}$
avec $E_a \ge
0$, ne peut que croître avec la température.
\textbf{6.} $-\dd[\ce{NO}]/\dd t = 2v$, $-\dd[\ce{O2}]/\dd t = v$.
\textbf{7.} \ce{2NO -> N2O2} plus \ce{N2O2 + O2 -> 2NO2} donne
\ce{2NO + O2 -> 2NO2}~; l'intermédiaire est \ce{N2O2}.
\textbf{8.} Tous deux bimoléculaires (et l'inverse du premier
monomoléculaire).
\textbf{9.} $v = k_2[\ce{N2O2}][\ce{O2}]$.
\textbf{10.} $[\ce{N2O2}] = K_1[\ce{NO}]^2$ avec $K_1 = k_1/k_{-1}$.
\textbf{11.} $v = k_2K_1[\ce{NO}]^2[\ce{O2}]$~: la loi observée.
\textbf{12.} $k = k_1k_2/k_{-1}$.
\textbf{13.} $\dd[\ce{N2O2}]/\dd t = k_1[\ce{NO}]^2 - k_{-1}[\ce{N2O2}] -
k_2[\ce{N2O2}][\ce{O2}]$.
\textbf{14.} En l'annulant~: $[\ce{N2O2}] = k_1[\ce{NO}]^2/(k_{-1} +
k_2[\ce{O2}])$.
\textbf{15.} $v = k_2[\ce{N2O2}][\ce{O2}]$ donne la loi annoncée.
\textbf{16.} $k_{-1} \gg k_2[\ce{O2}]$~: le dimère se dissocie bien plus
souvent qu'il ne rencontre le dioxygène, et le premier acte reste à
l'équilibre.
\textbf{17.} Si $k_2[\ce{O2}] \gg k_{-1}$, $v = k_1[\ce{NO}]^2$~: ordre 2
par rapport à \ce{NO} et 0 par rapport à \ce{O2}.
\textbf{18.} L'ordre 1 mesuré par rapport à \ce{O2}~: la première
limite.
\textbf{19.} $\ln k = \ln k_1 + \ln k_2 - \ln k_{-1}$.
\textbf{20.} En dérivant par rapport à $T$ et en utilisant
$\dd\ln k_i/\dd T
= E_i/RT^2$~: $E_a = E_1 + E_2 - E_{-1}$.
\textbf{21.} Le chauffage accélère un peu l'acte lent ($E_2$ faible) mais
déplace fortement le pré-équilibre vers \ce{NO} ($E_{-1}$ grande)~: moins
de dimères, et la vitesse nette diminue.
\textbf{22.} Le dimère se trouve $E_1 - E_{-1} = \qty{-38}{kJ/mol}$
au-dessous des réactifs et le second \omterm{def:b1:elementary-steps:energy-profile}{état de transition}
$E_2 = \qty{15}{kJ/mol}$ au-dessus du dimère~: \qty{23}{kJ/mol}
au-dessous des réactifs. Le point le plus haut du chemin est le premier
\omterm{def:b1:elementary-steps:energy-profile}{état de transition}, à seulement \qty{2}{kJ/mol}.
\textbf{23.} $\exp\big(\frac{23000}{8.314}(\frac{1}{400} - \frac{1}{300})\big)
= 0.10$~: dix fois plus lente à \qty{400}{K}.
\textbf{24.} $E_a = 2 - 40 + 15 = \textbf{\qty{-23}{kJ/mol}}$, la valeur
mesurée dans la partie I~: le mécanisme explique à la fois la \omterm{def:b1:rate-laws:order}{loi de
vitesse} et son étrange dépendance en température.
\end{solution}
